% Generated draft: requires mathematical and source review.
\documentclass[11pt]{article}
\usepackage[margin=25mm]{geometry}
\usepackage{fontspec}
\setmainfont{Latin Modern Roman}
\usepackage{amsmath,amssymb,mathtools,unicode-math}
\setmathfont{Latin Modern Math}
\usepackage{enumitem,needspace,xurl,hyperref}
\setlength{\parindent}{0pt}
\setlength{\parskip}{5pt}
\setlength{\emergencystretch}{3em}
\newcommand{\N}{\mathbb{N}}
\newcommand{\Z}{\mathbb{Z}}
\newcommand{\Q}{\mathbb{Q}}
\newcommand{\R}{\mathbb{R}}
\newcommand{\C}{\mathbb{C}}
\newcommand{\F}{\mathbb{F}}
\newcommand{\A}{\mathbb{A}}
\newcommand{\PP}{\mathbb{P}}
\newcommand{\E}{\mathbb{E}}
\newcommand{\eps}{\varepsilon}
\newcommand{\dd}{\,\mathrm{d}}
\begin{document}
\section*{UnsolvedMath ID: 30002374}
\textbf{Source problem code:} catalogue entry
\subsection{Definitions and mathematical statement}
Let \((M,\omega)\) be a symplectic manifold, and let \(\varphi\colon M\to M\) be a Hamiltonian diffeomorphism, meaning the time-one map of a Hamiltonian flow defined through time one. Define
\[
\operatorname{Fix}(\varphi)=\{x\in M:\varphi(x)=x\}.
\]
A periodic point is a point \(x\) for which \(\varphi^k(x)=x\) for some integer \(k\geq 1\). Its least period is the smallest such integer, and its periodic orbit is the finite set of its iterates. In the formulation below, distinct periodic orbits mean distinct such finite sets; counting different iterates of the same orbit separately would not express the intended infinitude assertion. The supplied record does not explicitly specify this counting convention or impose a contractibility condition.

\textbf{Hofer--Zehnder conjecture, as described in the supplied record.} A Hamiltonian diffeomorphism having more than the necessary number of fixed points has infinitely many distinct periodic orbits. To display the logical form, write \(N_{\mathrm{req}}(M,\omega)\) for the unspecified ``necessary number'' in the record. The proposed implication is
\[
\#\operatorname{Fix}(\varphi)>N_{\mathrm{req}}(M,\omega)
\quad\Longrightarrow\quad
\#\{\text{distinct periodic orbits of }\varphi\}=\infty.
\]
This is a schematic statement: the record does not define \(N_{\mathrm{req}}\), specify the admissible class of symplectic manifolds, or state the hypotheses under which that threshold is intended to apply. Consequently, the display is not a fully specified conjecture over all symplectic manifolds.

The record singles out the following \emph{expected case} on complex projective space:
\[
\varphi\in\operatorname{Ham}(\mathbb{CP}^{n},\omega),
\qquad
\#\operatorname{Fix}(\varphi)>n+1
\quad\Longrightarrow\quad
\#\{\text{distinct periodic orbits of }\varphi\}=\infty,
\]
where \(\operatorname{Ham}(M,\omega)\) denotes the Hamiltonian diffeomorphisms of \((M,\omega)\). The qualifier ``in the expected case'' is retained: the supplied record does not specify the symplectic form \(\omega\), the range of \(n\), or any further hypotheses intended for this implication.

The supplied background assessment invokes Franks's two-dimensional theorem and higher-dimensional results concerning hyperbolic fixed points as substantial special cases. It also describes an unresolved proposed extension from hyperbolic to general non-elliptic fixed points. No precise statement or hypotheses for that extension are supplied, so it cannot be identified here with an additional mathematical assertion.
\subsection{Short English statement}
The supplied Hofer--Zehnder conjecture predicts infinitely many distinct periodic orbits when a Hamiltonian diffeomorphism has more fixed points than a required minimum, with \(n+1\) suggested for \(\mathbb{CP}^{n}\) in an expected case. The general minimum, admissible manifolds, and qualifications for that expected case remain unspecified. The supplied partially solved label and literature assessment require independent review.
\subsection{Sources}
\begin{itemize}
\item Local registry: \texttt{data/unsolvedmath/problems.json}, supplied record with permanent UnsolvedMath ID \(30002374\), source problem code \texttt{catalogue entry}, title \emph{The Hofer--Zehnder Periodic-Orbit Conjecture}. The record attributes the conjecture by the name Hofer--Zehnder, lists proposed year 2013, and supplies no individual proposer.
\item Supplied source citation: \emph{Dynamische Systeme (2014)}, Oberwolfach Reports, pp.\ 1975--2033, DOI \texttt{10.4171/owr/2013/34}: \url{https://doi.org/10.4171/owr/2013/34}. The literature list also labels this source \emph{Oberwolfach Report 34/2013}. The discrepancy between the citation's 2014 label and its 2013 report identifiers is unresolved here.
\item Supplied literature source: \emph{Invariant sets and hyperbolic closed Reeb orbits}, \url{https://doi.org/10.1016/j.aim.2025.110709}.
\item Literature assessment dated exactly \texttt{2026-08-21}: supplied dated curation, not independently verified. It labels the problem partially solved, cites Franks's two-dimensional theorem and higher-dimensional hyperbolic-fixed-point results as substantial special cases, and describes the extension to general non-elliptic fixed points as unresolved. The background expressly requests independent verification of time-sensitive status and cited claims.
\item The registry labels the statement \texttt{corrected\_verified} and reports that its mathematical content was checked against the official Oberwolfach report, with notation, context, and question form normalized without changing the source-level claim. This is supplied verification metadata, not independent verification performed for this draft.
\end{itemize}
\end{document}
